Week 5: Understanding under threat --- quantum mechanics and AI

%% first part about QM, will take two thirds of the time, i.e. 2 x 45
%% minute lectures. Second part about AI will be 1 x 45 minute lecture


%% First part: Quantum Mechanics %%
%%%

\section{Introduction}

%% Today we will take two perspectives -- the second of which is not
%% typical in the way philosophers these days talk about foundations
%% of QM
%% 1. First-order perspective: thinking about which approach to QM we
%% think is the best
%% 2. Second-order perspective: conflict and consensus in science. Should we be
%% concerned that there is still no consensus about foundations  of
%% QM? 

%% I want to explain to the students that there was a foundational
%% crisis in physics between 1901 (Planck's discovery) and 1927. It
%% looked to be solved --- consolidation in Copenhagen around Bohr's
%% school.

%% But the crisis got revived in the 1960s, after the work of David
%% Bohm, Hugh Everett, and John Bell. Debate is raging now, not so
%% much at the cutting edges of physics --- where the theory works
%% well enough --- but among philosophers.

\item The early pioneers were acutely aware that they did \emph{not}
  understand what was going on
\item Bohr thought that our ability to understand the physical world
  was genuinely up for grabs.
\item His final analysis (circa 1927) was that understanding could be
  obtained, but changing what it means to ``understand''.

\section{Framework of quantum mechanics}

\begin{frame}

\item States as elements of a (complex) vector space. Illustrate with
  the Bloch sphere
  
\item Quantities as operators 

From the previous two items it follows that there is something a bit
puzzling going on with \emph{superpositions} of states

\end{frame}

\begin{frame}
  
[The uncertainty relation is basically Fourier analysis: the lower the
dispersion in one quantity, the higher in the other quantity. To
Claude: Illustrate with spin-$y$ and spin-$z$.]

\end{frame}

\begin{frame}{Composite systems}

  Composite systems are represented with the \textbf{tensor product
    formalism}, which is characterized by:

  \[ x\otimes (y+z) \: = \: (x\otimes y)+(x\otimes z) \]

  This formalism implies the existence of \textbf{entangled states},
  i.e.\ states that cannot be represented as a product of pure states
  for the two subsystems.

  For discussion (2 minutes): What features does a system have if it
  is entangled with another system?
  

\end{frame}  

\begin{frame}{Dynamics}

  The dynamics is \textbf{linear}:

  \[ U(x+y) \: = \: Ux + Uy \]

\end{frame}


\begin{frame}{Born rule}

  \[ \mathrm{Pr}_{\psi }(A=\lambda ) \: = \: \langle \psi ,E^A(\lambda
  )\psi \rangle  \]



\end{frame}



  QM is often taught as a recipe for computing the probabilities of
  measurement outcomes. That's a very \emph{antirealist} way of
  treating the theory. The problem that follows --- the measurement
  problem --- is supposed to make us worry about the consistency of
  maintaining such an antirealist position.

  Question for discussion: do we \emph{understand} QM if it's just a
  recipe for computing measurement outcomes?

\section{The measurement problem}

\item Definition of an ideal measurement of $A$
  -- states of $A$ are paired with states of the measuring device
  -- if the object is in state $x_i$ and the measuring device is in
  the ready state, then the composite object ends in state $x_i\otimes y_i$
\item Derivation of the problem
\item Collapse postulate and philosophical objections to it 
\item Historical aside: what happened between the 1930s and 1950s so
  that it became textbook QM with the collapse postulate
\item A little about the myth of the ``Copenhagen interpretation''
\item Bohr never suggested that the measurement problem is solved by
  collapse

\begin{frame}{Decoherence}

  But isn't the measurement problem solved by decoherence?

  Long story short: No.

  1. The reduced state is not uniquely interpretable as a mixture over
  the states of concern

  2. The state is not \emph{really} a mixture

\end{frame}
 

\section{Interpretations of QM}

\item Reject unitary dynamics. e.g. Ghirardi-Rimini-Weber
\item Hidden variables. e.g. Bohmian mechanics
\item Reject uniqueness of measurement outcomes. e.g. Everett
\item Interpretations that are ``less ontological'' in nature:
  relational, QBist, pragmatic (Healey)

\begin{frame}

  %% A slide or two to consider problems with the various
  %% interpretations
\item Bohm

  -- Technically inconsistent with STR

  -- Violates the spirit of the principle of conservation of momentum
  

\item GRW

  -- Unclear whether the wavefunction should be treated as some kind of
  matter density

  -- Wavefunction lives in high dimensional configuration space 

\item Everett

  -- Since the theory is deterministic, where do the probabilities come from?

  -- Preferred basis problem

  -- The no-hidden variables theorems imply that the worlds (branches) cannot be classical either

  

\end{frame}
  


%% now we shift over to quantum non-locality, passing through
%% complementarity, EPR, and no hidden variables theorems

\section{Complementarity}

%% here I'll give Bohr's famous picture of the two versions of the
%% single-slit experiment: one where the screen is bolted down, and
%% one where the screen freely floats on a infinitely loose spring

%% but I won't leave it there with pictures: I'll explain how he said
%% there is a complementarity between visualizability (spacetime
%% coordinatization) and definability (conservation of momentum)


\section{No hidden variables theorems}

%% remind them that Bohr didn't rest his argument on the no hidden
%% variables theorems. He knew about von Neumann's, but he didn't
%% think it was physically convincing

\item von Neumann's argument in Grundlagen der Quantenmechanik
\item Bell-Kochen-Specker


\section{Bell's theorem}

%% the focus here will be a bit different than the normal discussions
%% in analytic philosophy of physics. We aren't going to wring our
%% hands about locality and realism. We are more focused on whether
%% Bell's theorem shows limitations on our ability to "understand"
%% what is going on in the quantum world.

\item Derivation of an inequality from assumption of local hidden
  variables
\item Jarrett analysis: Realism plus locality implies inequality
\item Singlet state violates inequality [Hi Claude: choose spin
  operators at appropriate angles so that we can explicitly show
  computation of maximal violation of Bell's inequality. I think it
  achieves the value $2\sqrt{2}$?]
\item 
  

  \section{Conclusion}

  What can you do about the QM foundations ``crisis''?

  \item Some people have chosen to defect from mainstream physics
    because they think this issue is so important
  \item I don't think we should just tolerate it. QM is one thing ---
    the threat from AI is another thing. If we think it's OK to use QM
    without understanding, then why not just let AI do our science?

%% this is my transition to the AI part     


%%% Second part: AI     
  

%%% Local Variables:
%%% mode: latex
%%% TeX-master: t
%%% End:
